\chapter*{Заключение}
\addcontentsline{toc}{chapter}{Заключение}
\label{ch:conclusion}

В ходе лабораторной работы была развёрнута система автоматизации
\textbf{Ansible} на базе уже настроенных виртуальных машин из
лабораторных работ~5--8.

Выполненные задачи:
\begin{enumerate}
  \item На клиентских ВМ \texttt{desktop1} и \texttt{wordpress} установлены
        \texttt{openssh-server} и \texttt{python3} — обязательные зависимости
        для работы Ansible managed~node.
  \item На \texttt{gateway} установлен Ansible через \texttt{pip3},
        сгенерирован SSH-ключ \texttt{ed25519} без пароля.
  \item Публичный ключ распространён на оба клиента с помощью
        \texttt{ssh-copy-id}, что обеспечило беспарольный доступ.
  \item Созданы файлы \texttt{/etc/ansible/hosts} и \texttt{ansible.cfg};
        связность подтверждена командой \texttt{ansible clients -m ping}.
  \item Написан и запущен плейбук \texttt{monitoring.yml}, который собрал
        с каждого клиента: имя хоста, IP-адрес, версию ОС
        и объём свободного места на корневом разделе.
  \item Результаты сохранены в файлах
        \texttt{/etc/ansible/monitoring/*\_info.txt} на control~node.
\end{enumerate}

Использование \texttt{delegate\_to: localhost} позволило записывать
данные непосредственно на управляющий сервер без дополнительных сетевых
протоколов. Агентless-архитектура Ansible подтвердила своё преимущество:
клиентские ВМ потребовали минимальной настройки.

\endinput
